\section{模型评价与适用范围}\label{sec:evaluation}

本文没有用一个统一模型强行覆盖四个问题，而是按各问的数据结构选择对应方法。第一问
以组合内相邻观测保留局部变化，并用线性回归概括已测范围内的平均方向；第二问只在
共同温度和匹配组合之间比较，避免把样本构成变化误写成温度或组合影响；第三问把实测
推荐与缺测核查分开；第四问则针对前三问暴露的数据缺口配置有限实验。这样的分工使
每项结论都能回到相应的观测单位和比较范围。

本文采用的检验各有明确边界。独立复算由与主程序分离的计算路径重新读取输入并核对
公式、排序和程序输出，只验证计算一致性；第一问逐点
删除检查平均斜率方向是否受单点支配；第二问整组和整温度删除检查结论对比较对象与
温度范围的依赖；第三问内部遮点衡量缺测估计在已有数据结构中的预测误差；第四问逐项
删除只检查五次实验的职责是否存在冗余。这些检验都不能替代同条件重复、显著性检验或
全域误差估计。

模型的主要限制来自实验设计本身。多数“组合—温度”条件只有一次观测，不同组合的
温度覆盖也不完全相同，故共同温度之外的结果只能作匹配补充或个案说明。附件2又只有
一次时间实验，文献只能提供与现象相容的可能机制，不能证明具体失活路径。

因此，本文结论只适用于附件给出的21种组合、实测温度和观测时段。若要推广到连续温度
最优、新催化剂配方、长期运行、工业规模、安全性或经济性，需要增加重复实验并记录
批次与设备状态，再根据新增数据重新选择曲线形式和误差模型。
